% *======================================================================*
%  Cactus Thorn template for ThornGuide documentation
%  Author: Ian Kelley
%  Date: Sun Jun 02, 2002
%  $Header$
%
%  Thorn documentation in the latex file doc/documentation.tex
%  will be included in ThornGuides built with the Cactus make system.
%  The scripts employed by the make system automatically include
%  pages about variables, parameters and scheduling parsed from the
%  relevant thorn CCL files.
%
% *======================================================================*

% If you are using CVS use this line to give version information
% $Header$

\documentclass{article}

% Use the Cactus ThornGuide style file
% (Automatically used from Cactus distribution, if you have a
%  thorn without the Cactus Flesh download this from the Cactus
%  homepage at www.cactuscode.org)
\usepackage{../../../../doc/latex/cactus}

\begin{document}

% The author of the documentation
\author{Yuhua Li \textless liyhts@yeah.net\textgreater}

% The title of the document (not necessarily the name of the Thorn)
\title{Geodesic}

% the date your document was last changed
\date{September 9 2026}

\maketitle

% Do not delete next line
% START CACTUS THORNGUIDE

% Add all definitions used in this documentation here
%   \def\mydef etc

% Add an abstract for this thorn's documentation
\begin{abstract}
\texttt{Geodesic} is a pure geodesic particle tracker for Cactus and the
Einstein Toolkit.  It integrates the geodesic equation for an ensemble of
test particles in a fixed spacetime, using either the analytic Kerr--Schild
metric and Christoffel symbols built into the thorn (\texttt{Exact = yes})
or a metric read from the \texttt{ADMBase} grid, reconstructed by 27-point
local Lagrange polynomial interpolation (\texttt{Exact = no}).  Each
particle is integrated independently in proper time with the adaptive
GSL \texttt{rkf45} integrator; the 4-velocity is renormalized to
$u_\mu u^\mu = -1$ (future-directed) at every step; particles that
leave the domain are handled by a run-time selectable loss policy
(drop, flag, or respawn).  Because the interpolation is local to
each particle, the particles are embarrassingly parallel and are
advanced with OpenMP.
\end{abstract}

% The following sections are suggestive only.
% Remove them or add your own.

\section{Introduction}

\texttt{Geodesic} tracks test particles on a prescribed, fixed background
spacetime.  It is a pure geodesic integrator: the particles feel gravity
only, there is no electromagnetic or any other force, and the metric is
not evolved.  This makes the thorn useful as a controlled diagnostic
tool, for example to study how freely falling test particles move in a
Kerr background, to verify geodesic integrators against analytic
orbits, or to generate particle trajectories for visualization.

The thorn is deliberately minimal.  It runs inside a standard Cactus
build (Carpet driver, \texttt{CartGrid3D} grid), reads the metric either
analytically or from the grid, and advances one particle system per
Cactus iteration.

\section{Physical System}

The particles follow the geodesic equation
\begin{equation}
\frac{d^2 x^\mu}{d\tau^2} + \Gamma^\mu_{\alpha\beta}
\, u^\alpha u^\beta = 0,
\label{EinsteinAnalysis_Geodesic_geodesic_eq}
\end{equation}
where $\tau$ is the proper time and
$u^\mu = dx^\mu/d\tau$ is the 4-velocity, which satisfies the
normalization constraint
\begin{equation}
g_{\mu\nu} u^\mu u^\nu = -1,
\label{EinsteinAnalysis_Geodesic_norm}
\end{equation}
with the future-directed branch selected.  Equation
(\ref{EinsteinAnalysis_Geodesic_geodesic_eq}) is integrated as a first
order system for $(x^\mu, u^\mu)$.

The default background is the Kerr metric in Kerr--Schild Cartesian
coordinates, specified by the mass $M$ and spin $a$ (parameters
\texttt{M}, \texttt{a}).  The coordinates are the standard Cartesian
$(t, x, y, z)$ of the \texttt{CartGrid3D} grid, so particle positions
can be compared directly with grid coordinates.

\section{Numerical Implementation}

\subsection{Metric and Christoffel symbols}

Two modes are available, selected by the \texttt{Exact} parameter:

\begin{itemize}
\item \texttt{Exact = yes}: the metric components and all Christoffel
  symbols are evaluated from the analytic Kerr--Schild formulas built
  into the thorn.  No grid data is read, so the result is independent
  of the grid resolution; the grid is only needed by the driver.
\item \texttt{Exact = no}: the metric ($g_{ij}$, lapse $\alpha$, shift
  $\beta^i$) is read from the \texttt{ADMBase} grid, and the
  Christoffel symbols are assembled at each particle position by
  \emph{27-point local Lagrange interpolation}: a $3 \times 3 \times 3$
  stencil of grid values around the particle is fit by a local
  polynomial, of which the required spatial derivatives are taken
  analytically.  The time derivatives of the metric required by
  $\Gamma^\mu_{00}$ and $\Gamma^\mu_{0\nu}$ are computed with a
  2nd-order backward difference $(3 L_0 - 4 L_1 + L_2)/(2\,\Delta t)$
  over the three most recent \texttt{ADMBase} time levels, so the mode
  also works for time-dependent metrics.  A stationary metric is
  detected automatically (the two oldest levels agree to relative
  $10^{-12}$ on a subsample), in which case $dt(g) = 0$ exactly.
  The interpolation is local to each particle (no global polynomial is
  ever fit), so the cost per particle is independent of the particle
  number and the particles are fully independent.
\end{itemize}

\subsection{Time integration}

Each particle is integrated in proper time $\tau$ with the adaptive
GSL integrator \texttt{rkf45} (runge--kutta--fehlbert).  One Cactus
iteration advances the coordinate time by the fixed increment
\texttt{step} (or by \texttt{cctk\_delta\_time} if \texttt{step} = 0);
inside it, the integrator takes its own adaptive $\tau$ sub-steps until
the target coordinate time is reached.  With a time-dependent grid
metric the connection is frozen at the grid's current time within one
particle step, so when any refinement level's metric is not static the
\texttt{step} increment is capped at the grid time step of those
levels (a one-time warning is issued when the cap engages).

After each Cactus step the 4-velocity of every particle is
renormalized: $u^0$ is solved from the quadratic equation
(\ref{EinsteinAnalysis_Geodesic_norm}), and the future-directed root is
selected.  Inside the ergosphere the quadratic has two positive roots;
the thorn selects the one continuous with the previous step's value.
Particles that leave the domain, or for which no future-directed
solution exists, are handled by the \texttt{particle\_loss\_policy}
keyword: \texttt{drop} (the default) sets their state to NaN and stops
tracking them; \texttt{flag} freezes them at their last state;
\texttt{respawn} (legacy) re-places them at a random position inside a
safe box (outside the excision sphere, inside the initial shell) with
zero spatial velocity, and the next renormalization step then assigns
them a valid $u^0$.

\subsection{Parallelism and output}

Because both metric evaluations are local to each particle, the
particle loop is parallelized with OpenMP
(\texttt{\#pragma omp parallel for}) over the particle index.  The
particle arrays are stored with \texttt{DISTRIB=CONSTANT}
(replicated on every MPI rank); the integration runs on a single rank
(the other ranks idle during the particle step, and multi-process
runs are verified bit-identical against the single-rank result).
Multi-level (AMR) runs are supported for up to three Carpet
refinement levels (\texttt{geodesic\_amr\_levels}, default 3; unused
levels are skipped): on every refinement-level call the thorn
acquires that level's grid geometry and metric fields, and each
particle's 27-point stencil is drawn from the finest level whose grid
contains the particle in its safety region, so particles inside a
refined region are automatically sampled at the finer grid spacing.
Refined regions are assumed to consist of one centered patch, and
the run must be single-process.

The evolved state is kept in the \texttt{particle\_arrays} group
($\tau$, $t$, $x, y, z$, $u^t, u^x, u^y, u^z$), which is a regular
Cactus group: checkpointable, and readable by any output thorn.
With \texttt{particle\_dump\_every} $= k > 0$ the thorn additionally
writes a plain-text dump of the full particle state to
\texttt{particle\_dump\_file} (default
\texttt{geodesic\_particles.txt}, parent directories are created as
needed) every $k$ iterations, starting with the initial state.  This
dump is intended for quick visualization and for the testsuite.

\section{Using This Thorn}

\subsection{Obtaining This Thorn}

The thorn is part of the \texttt{EinsteinAnalysis} arrangement, in the
directory \texttt{arrangements/EinsteinAnalysis/Geodesic}.  It is
available on GitHub at
\url{https://github.com/wylyhzh/Geodesic}.

\subsection{A Minimal Run}

A complete minimal parameter file is
\texttt{par/example.par}: one particle on a prograde equatorial
circular orbit (radius 11.989) in a Kerr background with $M = 2$,
$a = 0.5$, \texttt{Exact = yes}, on a $50^3$ grid, 32 iterations.
The essential lines are
\begin{verbatim}
ActiveThorns = "Boundary CartGrid3D CoordBase SymBase Time"
ActiveThorns = "Carpet CarpetLib CarpetReduce"
ActiveThorns = "IOUtil"
ActiveThorns = "ADMBase KerrSchild"
ActiveThorns = "Geodesic"
...
Geodesic::M          = 2.0
Geodesic::a          = 0.50
Geodesic::Exact      = "yes"
Geodesic::step       = 100.0
Geodesic::particle_n_total = 1
Geodesic::particle_n       = 1
Geodesic::particle_x[0] = 11.989
Geodesic::particle_tv[0] = 1.392812583392655
Geodesic::particle_yv[0] = 0.5600648870761632
Geodesic::particle_dump_every = 1
\end{verbatim}
The \texttt{particle\_tv/xv/yv/zv} parameters are the contravariant
4-velocity components $u^\mu = (u^t, u^x, u^y, u^z)$ in Kerr--Schild
coordinates (not a 3-velocity); the thorn checks
$g_{\mu\nu}u^\mu u^\nu = -1$ at initialization and uses the given
values if they are normalized.

\subsection{Tests}

The thorn ships a testsuite in \texttt{test/} with three tests, all
integrating the same prograde equatorial circular orbit
($r = 11.989$, $M = 2$, $a = 0.5$, \texttt{Exact = yes}, 32
iterations of 100 coordinate-time units):

\begin{itemize}
\item \texttt{circ\_orbit} compares the full particle state dump
  against reference data row by row; any change of the trajectory
  fails the test.
\item \texttt{circ\_orbit\_conserved} runs the dump through the
  postprocessor \texttt{util/check\_conserved}, which prints the two
  Kerr constants of motion $E$ and $L_z$ for every row, and compares
  them against the reference; a broken integrator shows up as a drift
  of $E$ / $L_z$.  On a correct run both are constant to about nine
  digits over the whole orbit.
\item \texttt{circ\_orbit\_mpi} is the \texttt{circ\_orbit} run with
  2 MPI ranks (requires \texttt{mpirun}); the integration runs on a
  single rank, so the dump must be bit-identical to the single-rank
  reference.
\end{itemize}

Run the testsuite with
\begin{verbatim}
cd <cactus>
CCTK_TESTSUITE_RUN_TESTS="Geodesic/circ_orbit Geodesic/circ_orbit_conserved \
Geodesic/circ_orbit_mpi" gmake sim-testsuite
\end{verbatim}

\subsection{Interaction With Other Thorns}

\texttt{Geodesic} requires \texttt{Carpet} (the only thorn in this
build providing Cactus group storage, which is needed by the
\texttt{particle\_arrays} variable) and auto-activates \texttt{GSL}.
In \texttt{Exact = no} mode the metric is read from the
\texttt{ADMBase} grid (the \texttt{schedule.ccl} declares the
corresponding \texttt{READS} dependencies), so any thorn that
provides the \texttt{ADMBase} variables can be used as the metric
source, with \texttt{KerrSchild} the standard choice.  No
hydrodynamics or other evolution thorn is needed; in
\texttt{Exact = yes} mode the grid content is not used at all.

% Do not delete next line
% END CACTUS THORNGUIDE

\end{document}
